\section{AHSME 1996}
        \begin{enumerate}
        	\item (\textit{AHSME 1996}) The addition below is incorrect. What is the largest digit that can be changed to make the addition correct?


 $\begin{tabular}{rr}&\ \texttt{6 4 1}\\ &\texttt{8 5 2}\\ +&\texttt{9 7 3}\\ \hline  &\texttt{2 4 5 6}\end{tabular}$


 

$\text{(A)}\ 4\qquad\text{(B)}\ 5\qquad\text{(C)}\ 6\qquad\text{(D)}\ 7\qquad\text{(E)}\ 8$


 

	\item (\textit{AHSME 1996}) Each day Walter gets $3$ dollars for doing his chores or $5$ dollars for doing them exceptionally well. After $10$ days of doing his chores daily, Walter has received a total of $36$ dollars. On how many days did Walter do them exceptionally well?


 $\text{(A)}\ 3\qquad\text{(B)}\ 4\qquad\text{(C)}\ 5\qquad\text{(D)}\ 6\qquad\text{(E)}\ 7$


 

	\item (\textit{AHSME 1996}) $\frac{(3!)!}{3!}=$


 $\text{(A)}\ 1\qquad\text{(B)}\ 2\qquad\text{(C)}\ 6\qquad\text{(D)}\ 40\qquad\text{(E)}\ 120$


 

	\item (\textit{AHSME 1996}) Six numbers from a list of nine integers are $7,8,3,5, 9$ and $5$. The largest possible value of the median of all nine numbers in this list is


 $\text{(A)}\ 5\qquad\text{(B)}\ 6\qquad\text{(C)}\ 7\qquad\text{(D)}\ 8\qquad\text{(E)}\ 9$


 

	\item (\textit{AHSME 1996}) Given that $0 < a < b < c < d$, which of the following is the largest?


 $\text{(A)}\  \frac{a+b}{c+d} \qquad\text{(B)}\ \frac{a+d}{b+c} \qquad\text{(C)}\  \frac{b+c}{a+d} \qquad\text{(D)}\  \frac{b+d}{a+c} \qquad\text{(E)}\ \frac{c+d}{a+b}$


 

	\item (\textit{AHSME 1996}) If $f(x) = x^{(x+1)}(x+2)^{(x+3)}$, then $f(0)+f(-1)+f(-2)+f(-3) =$


 $\text{(A)}\ -\frac{8}{9}\qquad\text{(B)}\ 0\qquad\text{(C)}\ \frac{8}{9}\qquad\text{(D)}\ 1\qquad\text{(E)}\ \frac{10}{9}$


 

	\item (\textit{AHSME 1996}) A father takes his twins and a younger child out to dinner on the twins' birthday. The restaurant charges $4.95$ for the father and $0.45$ for each year of a child's age, where age is defined as the age at the most recent birthday. If the bill is $9.45$, which of the following could be the age of the youngest child? 


 $\text{(A)}\ 1\qquad\text{(B)}\ 2\qquad\text{(C)}\ 3\qquad\text{(D)}\ 4\qquad\text{(E)}\ 5$


 

	\item (\textit{AHSME 1996}) If $3 = k\cdot 2^r$ and $15 = k\cdot 4^r$, then $r =$


 $\text{(A)}\ -\log_{2}5\qquad\text{(B)}\ \log_{5}2\qquad\text{(C)}\ \log_{10}5\qquad\text{(D)}\ \log_{2}5\qquad\text{(E)}\ \frac{5}{2}$


 

	\item (\textit{AHSME 1996}) Triangle $PAB$ and square $ABCD$ are in perpendicular planes. Given that $PA = 3, PB = 4$ and $AB = 5$, what is $PD$? 

                    \begin{figure}[H]
                        \centering
                        \includegraphics[width=0.5\textwidth]{Problem/AHSME_I_1996_9_Problem_Figure_1.png}
                    \end{figure}
                    


 $\text{(A)}\ 5\qquad\text{(B)}\ \sqrt{34} \qquad\text{(C)}\ \sqrt{41}\qquad\text{(D)}\ 2\sqrt{13}\qquad\text{(E)}\ 8$


 

	\item (\textit{AHSME 1996}) How many line segments have both their endpoints located at the vertices of a given cube?


 $\text{(A)}\ 12\qquad\text{(B)}\ 15\qquad\text{(C)}\ 24\qquad\text{(D)}\ 28\qquad\text{(E)}\ 56$


 

	\item (\textit{AHSME 1996}) Given a circle of radius $2$, there are many line segments of length $2$ that are tangent to the circle at their midpoints. Find the area of the region consisting of all such line segments.


 $\text{(A)}\ \frac{\pi} 4\qquad\text{(B)}\ 4-\pi\qquad\text{(C)}\ \frac{\pi} 2\qquad\text{(D)}\ \pi\qquad\text{(E)}\ 2\pi$


 

	\item (\textit{AHSME 1996}) A function $f$ from the integers to the integers is defined as follows: 


 
 \[f(x) =\begin{cases}n+3 &\text{if n is odd}\\ n/2 &\text{if n is even}\end{cases}\]


 Suppose $k$ is odd and $f(f(f(k))) = 27$.  What is the sum of the digits of $k$?


 $\textbf{(A)}\ 3\qquad\textbf{(B)}\ 6\qquad\textbf{(C)}\ 9\qquad\textbf{(D)}\ 12\qquad\textbf{(E)}\ 15$


 

	\item (\textit{AHSME 1996}) Sunny runs at a steady rate, and Moonbeam runs $m$ times as fast, where $m$ is a number greater than 1.  If Moonbeam gives Sunny a head start of $h$ meters, how many meters must Moonbeam run to overtake Sunny? 


 $\text{(A)}\ hm\qquad\text{(B)}\ \frac{h}{h+m}\qquad\text{(C)}\ \frac{h}{m-1}\qquad\text{(D)}\ \frac{hm}{m-1}\qquad\text{(E)}\ \frac{h+m}{m-1}$


 

	\item (\textit{AHSME 1996}) Let $E(n)$ denote the sum of the even digits of $n$. For example, $E(5681) = 6+8 = 14$.  Find $E(1)+E(2)+E(3)+\cdots+E(100)$


 $\text{(A)}\ 200\qquad\text{(B)}\ 360\qquad\text{(C)}\ 400\qquad\text{(D)}\ 900\qquad\text{(E)}\ 2250$


 

	\item (\textit{AHSME 1996}) Two opposite sides of a rectangle are each divided into $n$ congruent segments, and the endpoints of one segment are joined to the center to form triangle $A$. The other sides are each divided into $m$ congruent segments, and the endpoints of one of these segments are joined to the center to form triangle $B$. [See figure for $n=5, m=7$.] What is the ratio of the area of triangle $A$ to the area of triangle $B$?


 
                    \begin{figure}[H]
                        \centering
                        \includegraphics[width=0.5\textwidth]{Problem/AHSME_I_1996_15_Problem_Figure_1.png}
                    \end{figure}
                    


 $\text{(A)}\ 1\qquad\text{(B)}\ m/n\qquad\text{(C)}\ n/m\qquad\text{(D)}\ 2m/n\qquad\text{(E)}\ 2n/m$


 

	\item (\textit{AHSME 1996}) A fair standard six-sided dice is tossed three times. Given that the sum of the first two tosses equal the third, what is the probability that at least one "2" is tossed?


 $\text{(A)}\ \frac{1}{6}\qquad\text{(B)}\ \frac{91}{216}\qquad\text{(C)}\ \frac{1}{2}\qquad\text{(D)}\ \frac{8}{15}\qquad\text{(E)}\ \frac{7}{12}$


 

	\item (\textit{AHSME 1996}) In rectangle $ABCD$, angle $C$ is trisected by $\overline{CF}$ and $\overline{CE}$, where $E$ is on $\overline{AB}$, $F$ is on $\overline{AD}$, $BE=6$ and $AF=2$. Which of the following is closest to the area of the rectangle $ABCD$?

                    \begin{figure}[H]
                        \centering
                        \includegraphics[width=0.5\textwidth]{Problem/AHSME_I_1996_17_Problem_Figure_1.png}
                    \end{figure}
                    
$\text{(A)}\ 110\qquad\text{(B)}\ 120\qquad\text{(C)}\ 130\qquad\text{(D)}\ 140\qquad\text{(E)}\ 150$


 

	\item (\textit{AHSME 1996}) A circle of radius $2$ has center at $(2,0)$. A circle of radius $1$ has center at $(5,0)$. A line is tangent to the two circles at points in the first quadrant. Which of the following is closest to the $y$-intercept of the line?


 $\text{(A)}\ \sqrt{2}/4\qquad\text{(B)}\ 8/3\qquad\text{(C)}\ 1+\sqrt 3\qquad\text{(D)}\ 2\sqrt 2\qquad\text{(E)}\ 3$


 

	\item (\textit{AHSME 1996}) The midpoints of the sides of a regular hexagon $ABCDEF$ are joined to form a smaller hexagon. What fraction of the area of $ABCDEF$ is enclosed by the smaller hexagon? 


 
                    \begin{figure}[H]
                        \centering
                        \includegraphics[width=0.5\textwidth]{Problem/AHSME_I_1996_19_Problem_Figure_1.png}
                    \end{figure}
                    


 $\text{(A)}\ \frac{1}{2}\qquad\text{(B)}\ \frac{\sqrt 3}{3}\qquad\text{(C)}\ \frac{2}{3}\qquad\text{(D)}\ \frac{3}{4}\qquad\text{(E)}\ \frac{\sqrt 3}{2}$


 

	\item (\textit{AHSME 1996}) In the xy-plane, what is the length of the shortest path from $(0,0)$ to $(12,16)$ that does not go inside the circle $(x-6)^{2}+(y-8)^{2}= 25$?


 $\text{(A)}\ 10\sqrt 3\qquad\text{(B)}\ 10\sqrt 5\qquad\text{(C)}\ 10\sqrt 3+\frac{ 5\pi}{3}\qquad\text{(D)}\ 40\frac{\sqrt{3}}{3}\qquad\text{(E)}\ 10+5\pi$


 

	\item (\textit{AHSME 1996}) Triangles $ABC$ and $ABD$ are isosceles with $AB=AC=BD$, and $BD$ intersects $AC$ at $E$. If $BD$ is perpendicular to $AC$, then $\angle C+\angle D$ is


 
                    \begin{figure}[H]
                        \centering
                        \includegraphics[width=0.5\textwidth]{Problem/AHSME_I_1996_21_Problem_Figure_1.png}
                    \end{figure}
                    


 $\text{(A)}\ 115^\circ\qquad\text{(B)}\ 120^\circ\qquad\text{(C)}\ 130^\circ\qquad\text{(D)}\ 135^\circ\qquad\text{(E)}\ \text{not uniquely determined}$


 

	\item (\textit{AHSME 1996}) Four distinct points, $A$, $B$, $C$, and $D$, are to be selected from $1996$ points 
evenly spaced around a circle. All quadruples are equally likely to be chosen. 
What is the probability that the chord $AB$ intersects the chord $CD$? 


 $\text{(A)}\ \frac{1}{4}\qquad\text{(B)}\ \frac{1}{3}\qquad\text{(C)}\ \frac{1}{2}\qquad\text{(D)}\ \frac{2}{3}\qquad\text{(E)}\ \frac{3}{4}$


 

	\item (\textit{AHSME 1996}) The sum of the lengths of the twelve edges of a rectangular box is $140$, and 
the distance from one corner of the box to the farthest corner is $21$. The total 
surface area of the box is 


 $\text{(A)}\ 776\qquad\text{(B)}\ 784\qquad\text{(C)}\ 798\qquad\text{(D)}\ 800\qquad\text{(E)}\ 812$


 

	\item (\textit{AHSME 1996}) The sequence $1, 2, 1, 2, 2, 1, 2, 2, 2, 1, 2, 2, 2, 2, 1, 2, 2, 2, 2, 2, 1, 2,\ldots$ consists of $1$’s separated by blocks of $2$’s with $n$ $2$’s in the $n^{th}$ block. The sum of the first $1234$ terms of this sequence is


 $\text{(A)}\ 1996\qquad\text{(B)}\ 2419\qquad\text{(C)}\ 2429\qquad\text{(D)}\ 2439\qquad\text{(E)}\ 2449$


 

	\item (\textit{AHSME 1996}) Given that $x^2 + y^2 = 14x + 6y + 6$, what is the largest possible value that $3x + 4y$ can have? 


 $\text{(A)}\ 72\qquad\text{(B)}\ 73\qquad\text{(C)}\ 74\qquad\text{(D)}\ 75\qquad\text{(E)}\ 76$


 

	\item (\textit{AHSME 1996}) An urn contains marbles of four colors: red, white, blue, and green. When four marbles are drawn without replacement, the following events are equally likely: 


 (a) the selection of four red marbles; 


 (b) the selection of one white and three red marbles; 


 (c) the selection of one white, one blue, and two red marbles; and 


 (d) the selection of one marble of each color. 


 What is the smallest number of marbles satisfying the given condition? 


 $\text{(A)}\ 19\qquad\text{(B)}\ 21\qquad\text{(C)}\ 46\qquad\text{(D)}\ 69\qquad\text{(E)}\ \text{more than 69}$


 

	\item (\textit{AHSME 1996}) Consider two solid spherical balls, one centered at $(0, 0,\frac{21}{2})$ with radius $6$, and the other centered at $(0, 0, 1)$ with radius $\frac{9}{2}$. How many points with only integer coordinates (lattice points) are there in the intersection of the balls? 


 $\text{(A)}\ 7\qquad\text{(B)}\ 9\qquad\text{(C)}\ 11\qquad\text{(D)}\ 13\qquad\text{(E)}\ 15$


 

	\item (\textit{AHSME 1996}) On a $4\times 4\times 3$ rectangular parallelepiped, vertices $A$, $B$, and $C$ are adjacent to vertex $D$. The perpendicular distance from $D$ to the plane containing
$A$, $B$, and $C$ is closest to 


 
                    \begin{figure}[H]
                        \centering
                        \includegraphics[width=0.5\textwidth]{Problem/AHSME_I_1996_28_Problem_Figure_1.png}
                    \end{figure}
                    


 $\text{(A)}\ 1.6\qquad\text{(B)}\ 1.9\qquad\text{(C)}\ 2.1\qquad\text{(D)}\ 2.7\qquad\text{(E)}\ 2.9$


 

	\item (\textit{AHSME 1996}) If $n$ is a positive integer such that $2n$ has $28$ positive divisors and $3n$ has $30$ positive divisors, then how many positive divisors does $6n$ have?


 $\text{(A)}\ 32\qquad\text{(B)}\ 34\qquad\text{(C)}\ 35\qquad\text{(D)}\ 36\qquad\text{(E)}\ 38$


 

	\item (\textit{AHSME 1996}) A hexagon inscribed in a circle has three consecutive sides each of length $3$ and three consecutive sides each of length $5$. The chord of the circle that divides the hexagon into two trapezoids, one with three sides each of length $3$ and the other with three sides each of length $5$, has length equal to $\frac{m}{n}$, where $m$ and $n$ are relatively prime positive integers. Find $m+n$.


 $\text{(A)}\ 309\qquad\text{(B)}\ 349\qquad\text{(C)}\ 369\qquad\text{(D)}\ 389\qquad\text{(E)}\ 409$


 

\end{enumerate}